\begin{abstract}
We study unsupervised human activity recognition on the UCI smartphone HAR dataset using finite-state Gaussian hidden Markov models. Following our proposal, we compare the same Gaussian HMM under two inference strategies: classical maximum-likelihood estimation with EM and Bayesian inference with Gibbs sampling. We assume a fully unsupervised setting with no activity labels used during training. To keep Bayesian inference tractable, we use weakly informative conjugate priors: Dirichlet priors for the simplex-constrained initial and transition probabilities, and Normal-Inverse-Wishart priors for Gaussian emission parameters. A toy-data sanity check based on cached Step 2 outputs verifies that exact Forward-Backward inference and the Bayesian sampler recover the planted hidden-state structure. On the HAR benchmark, we reduce the 561 engineered features to eight principal components, treat each subject as a separate sequence, and quantify robustness by held-out log-likelihood plus clustering degradation under added Gaussian noise. The resulting comparison shows that EM slightly outperforms the Bayesian model on clustering metrics, while the Bayesian model attains slightly better held-out log-likelihood and provides posterior uncertainty over state persistence.
\end{abstract}
